\rok{Март 2026.}{B}

Трећи поправни тест из Алгоритама и структура података

Први практични тест одржан 06.03.2026.

\zadatak{1}{Merge Sort}
Написати методу \texttt{merge} Merge Sort алгоритма за сортирање целобројног низа дужине $n$.

\textbf{Додатне напомене за решавање задатка:}
\begin{itemize}
	\item Улаз је несортиран низ целих бројева.
	\item Излаз је сортирани низ целих бројева.
	\item У template-у се налази класа \texttt{RandomArrayGenerator} коју треба искористити за генерисање низа случајних бројева.
	\item У оквиру класе \texttt{MergeSort} написати имплементацију за методу:
	\begin{Verbatim}[fontsize=\small]
private static void merge(long[] arr, long[] tmpArray, int lowerIndex, int middleIndex, int upperIndex)
	\end{Verbatim}
	\item Обезбедити код у Main методи за тестирање алгоритма.
\end{itemize}

\zadatak{2}{Замена места првом и последњем парном чвору}
Написати методу која у једном пролазу идентификује први и последњи парни чвор у листи и замени им њихове вредности (value). Остатак листе и везе између чворова морају остати нетакнути.

\textbf{Додатни захтеви за решавање задатка:}
\begin{itemize}
	\item Задатак решити у линеарном времену $O(n)$.
	\item Ако у листи има мање од два парна броја, метода не ради ништа.
	\item Захтеван потпис методе:
	\begin{Verbatim}[fontsize=\small]
public void swapEvenBounds()
	\end{Verbatim}
\end{itemize}

\textbf{Пример:}
\begin{itemize}
	\item \textbf{Улаз:} \texttt{1 -> 2 -> 3 -> 10}
	\item \textbf{Излаз:} \texttt{1 -> 10 -> 3 -> 2}
\end{itemize}

\zadatak{3}{Контрола протока у уском грлу (Gateway Traffic Controller)}
Мрежни пролаз (Gateway) прима пакете података. Постоје два ограничења: укупна количина података у бајтовима коју пролаз може да обради у секунди, и максималан број пакета (slotova) који могу истовремено бити у обради због ограниченог броја хардверских канала.

Опис задатка: Написати методу која симулира пролаз пакета кроз gateway. Метода прима низ \texttt{packetSizes}, максимални капацитет бајтова \texttt{maxBytes} и максималан број slotova \texttt{maxSlots}.

\begin{enumerate}
	\item Пакети пристижу један по један.
	\item Да би нови пакет ушао у обраду, мора бити испуњено:
	\begin{itemize}
		\item Тренутна сума бајтова у обради + нови пакет $\leq$ \texttt{maxBytes}.
		\item Тренутни број пакета у обради + 1 $\leq$ \texttt{maxSlots}.
	\end{itemize}
	\item Ако услови нису испуњени, gateway избацује најстарије пакете све док оба услова не буду задовољена.
	\item Ако је пакет сам по себи већи од \texttt{maxBytes}, он се игнорише.
\end{enumerate}

\textbf{Додатни захтеви за решавање задатка:}
\begin{itemize}
	\item Повратна вредност: \texttt{int[]} (низ величина пакета који су остали у обради на крају симулације).
	\item Пomoћна структура: користити достављену класу \texttt{Queue}.
	\item Потпис:
	\begin{Verbatim}[fontsize=\small]
public int[] simulateGateway(int[] packetSizes, int maxBytes, int maxSlots)
	\end{Verbatim}
\end{itemize}

\textbf{Пример:}
\begin{itemize}
	\item \textbf{Улаз 1:} \texttt{packetSizes = [10, 10, 10, 10, 50, 5], maxBytes = 60, maxSlots = 3}
	\item \textbf{Симулација корак по корак:}
	\begin{enumerate}
		\item Стиже \texttt{10}: Slotovi \texttt{(1/3)}, тежина \texttt{(10/60)}. Ok.
		\item Стиже \texttt{10}: Slotovi \texttt{(2/3)}, тежина \texttt{(20/60)}. Ok.
		\item Стиже \texttt{10}: Slotovi \texttt{(3/3)}, тежина \texttt{(30/60)}. Ok.
		\item Стиже \texttt{10}: Problem! Slotovi би били \texttt{4/3}. Избацује се најстарији \texttt{10}. Стање: \texttt{[10, 10, 10]}.
		\item Стиже \texttt{50}: Problem! Тежина би била \texttt{30+50=80}.
		\begin{itemize}
			\item Избацује се најстарији \texttt{10}. Тежина \texttt{20+50=70} (и даље превише).
			\item Избацује се следећи најстарији \texttt{10}. Тежина \texttt{10+50=60}. (Сада је OK).
			\item Стање: \texttt{[10, 50]}.
		\end{itemize}
		\item Стиже \texttt{5}: Slotovi су \texttt{2/3}, али тежина би била \texttt{60+5=65}.
		\begin{itemize}
			\item Избацује се најстарији \texttt{10}. Тежина \texttt{50+5=55}. (Сада је OK).
			\item Стање: \texttt{[50, 5]}.
		\end{itemize}
	\end{enumerate}
	\item \textbf{Излаз 1:} \texttt{[50, 5]}
\end{itemize}

\sablon{Шаблон трећег поправног теста март 2026. група Б}{2025/Mart/GrupaB/AISP2025_PopravniTest3_Test1_GrupaB.linq}
